\begin{abstract}
%%%%%%%%%%%%%%%%%%%%%%%%%%%%%%%%%%%%%%%%%%%
%\guideabs % delete me when you know what goes in this section
%%%%%%%%%%%%%%%%%%%%%%%%%%%%%%%%%%%%%%%%%%%
%% example filler text:
Road traffic accidents remain a persistent public safety concern, yet local authorities often lack transparent, evidence-based tools for identifying where risk concentrates on urban road networks. This study identifies and profiles road-accident hotspots in Stavanger, Norway, from 2010 to 2025 by integrating open Norwegian data sources: police-reported accident records from Felles Datakatalog, road characteristics from the Nasjonal vegdatabank (NVDB), and traffic volume from Statens vegvesen traffic registration points. We join accident locations to road sections in a GIS environment and detect hotspots using kernel density estimation (KDE) and Getis-Ord Gi*. We examine hotspot stability across methods, parameter settings, time periods, and severity levels, including killed or seriously injured (KSI) outcomes. Accident counts are then adjusted for traffic exposure to test whether high-volume corridors explain apparent concentrations or whether some sections carry disproportionate risk. Finally, we train an interpretable model based on traffic volume, speed limit, and section length on 70 percent of the road-section data and evaluate it on the remaining 30 percent, comparing its ranking of sections by relative accident risk against a baseline that ranks sections by past accident counts alone. Findings are communicated through an interactive web dashboard designed for planners and non-specialist readers. The study is observational and identifies statistical associations rather than causes; its principal limitations are the geographic precision of historical accident records and coverage gaps in traffic sensors. The expected contribution is an applied, reproducible workflow combining established spatial methods, exposure-adjusted risk profiling and transparent visualisation for a mid-sized Norwegian city.
\keywords{accident hotspot detection; geospatial analysis; kernel density estimation; risk profiling; Stavanger} %limit these to 3-5 terms
\end{abstract}